\begin{table}[htbp]
\centering
\caption{Summary of Major Union Army Pension Laws}
\label{tab:pension_laws}
\begin{tabular}{lp{3cm}p{4cm}p{4cm}}
\toprule
  Law & Date & Eligibility & Benefits \\
\midrule
  General Law & July 14, 1862 & Service-connected disability & \$8--\$30/mo by disability rating \\
  Dependent Pension Act & June 27, 1890 & Any disability (not service-connected) & \$6--\$12/mo \\
  \textbf{Service and Age Act} & \textbf{Feb 6, 1907} & \textbf{Age 62+, 90+ days service} & \textbf{\$12--\$20/mo by age} \\
  Sherwood Act$^{\dagger}$ & May 11, 1912 & Age 62+, liberalized & \$13--\$30/mo \\
\bottomrule
\end{tabular}
\begin{tablenotes}[flushleft]
\small
\item \textit{Notes:} Principal Union Army pension laws. The 1907 Act created the first age-based eligibility threshold at 62, providing \$12/month (ages 62--69), \$15/month (ages 70--74), and \$20/month (ages 75+) regardless of disability status. $^{\dagger}$The Sherwood Act postdates the 1910 census and does not affect the analysis; it is included for historical completeness.
\end{tablenotes}
\end{table}
